\documentclass[11pt,twocolumn]{article}

\usepackage[margin=0.75in]{geometry}
\usepackage{amsmath,amssymb}
\usepackage{booktabs}
\usepackage{hyperref}
\usepackage{xcolor}
\usepackage{enumitem}
\usepackage{caption}
\usepackage{tabularx}
\usepackage{array}
\usepackage{multirow}
\usepackage[T1]{fontenc}
\usepackage{lmodern}
\usepackage{microtype}
\usepackage{fancyhdr}
\usepackage{titlesec}

\hypersetup{colorlinks=true,linkcolor=blue,citecolor=blue,urlcolor=blue}

\titleformat{\section}{\large\bfseries}{\thesection.}{0.5em}{}
\titleformat{\subsection}{\normalsize\bfseries}{\thesubsection}{0.5em}{}

\pagestyle{fancy}
\fancyhf{}
\fancyfoot[C]{\thepage}
\renewcommand{\headrulewidth}{0pt}

\title{\textbf{Tier-Based Adaptive Tool Routing for\\Capability-Heterogeneous AI Agents}}
\author{%
  \textbf{Pranab Sarkar}\\
  Yantrikos\\
  \texttt{developer@pranab.co.in}\\
  \url{https://github.com/spranab}\\[0.5em]
  \small ORCID: 0009-0009-8683-1481\\[1em]
  \small\textbf{Project:} \url{https://github.com/yantrikos/tier}
}
\date{Version 1.0.0 --- March 2026}

\begin{document}
\twocolumn[
  \maketitle
  \begin{center}
    \rule{0.5\textwidth}{0.4pt}
  \end{center}
  \vspace{0.5em}
  \begin{abstract}
  \noindent Tool-using AI agents are deployed across models ranging from 1.5B to 35B+ parameters, yet frameworks present identical tool interfaces regardless of model capability. We evaluate seven adaptive presentation strategies across four models (1.5B--35B) with 80 tools using native tool calling APIs. A novel decomposition reveals that tool selection accuracy factors into \emph{family routing accuracy} and \emph{within-family selection accuracy}, where the latter is consistently high (89--98\%) across all model sizes. This means the tool selection problem is primarily a \textbf{family routing problem}. Our best strategies improve accuracy by +8--10pp while reducing prompt tokens by 83--92\%: hybrid presentation (8~detailed + 72~name-only tools) achieves 60\% for 1.5B models (+10pp), while semantic reordering with category hints achieves 88\% for 20B models (+8pp), matching 35B baseline performance. We validate these findings against a production system (YantrikOS) with 116+ tools, and release an open-source benchmark, SDK, and OpenClaw plugin. Benchmark data: 1,000+ inference calls across three experimental rounds.

  \medskip
  \noindent\textbf{Keywords:} tool use, function calling, small language models, adaptive routing, AI agents, model capability, YantrikDB
  \end{abstract}
  \vspace{1em}
]

% ======================================================================
\section{Introduction}

\subsection{The Problem}

Tool-using AI agents invoke external functions --- reading files, querying databases, sending messages, executing code. Frameworks such as OpenAI function calling, Anthropic tool use, and open-source platforms like OpenClaw present tools through structured schemas that models select from.

A critical assumption pervades these frameworks: \textbf{all models receive the same tool presentation}. Whether the model has 1.5B or 35B parameters, it receives identical descriptions, parameter schemas, and selection interfaces. With 80 tools, this produces 2,100--5,300 prompt tokens of tool metadata --- consuming 30--50\% of a small model's context window.

The consequences are measurable: a 1.5B model achieves only 50\% tool selection accuracy with 80 native tools, while a 35B model achieves 88\%. The gap is not in the model's ability to \emph{use} tools, but in its ability to \emph{navigate} a large tool space.

\subsection{Key Insight: The Decomposition}

We decompose tool selection into two stages:

\begin{equation}
P(\text{correct tool}) = P(\text{correct family}) \times P(\text{correct tool} \mid \text{correct family})
\label{eq:decomposition}
\end{equation}

Our experiments reveal that even a 1.5B model achieves \textbf{89\% within-family accuracy} --- comparable to the 98\% achieved by a 35B model. The bottleneck is family routing, not tool selection. This insight reframes the problem: instead of making small models handle more tools, we should help them find the right tool \emph{neighborhood}.

\subsection{Contributions}

\begin{enumerate}[leftmargin=*,nosep]
  \item A \textbf{decomposition of tool routing} into family selection and within-family disambiguation, showing the bottleneck is routing (56\% for 1.5B) not selection (89\% given correct family).
  \item An \textbf{empirical evaluation of seven strategies} across four model sizes using native tool calling APIs, with 1,000+ total inference calls.
  \item \textbf{Two practical strategies} that improve accuracy without additional LLM calls: hybrid presentation (+10pp for 1.5B) and semantic reordering with hints (+8pp for 20B).
  \item A \textbf{production-validated framework} in YantrikOS with 116+ tools, and an open-source SDK for tier-aware tool development.
\end{enumerate}

% ======================================================================
\section{Related Work}

\textbf{Tool use in LLMs.} Function calling is now standard across OpenAI~\cite{openai_fc}, Anthropic, and open-source models (Qwen~\cite{qwen25}, Llama~\cite{llama3}). These systems present all available tools regardless of model capability.

\textbf{Small models for agents.} AgentFlux~\cite{agentflux} decouples tool selection from argument generation, improving 7B accuracy by 46\% but still presenting all tools. TinyLLM~\cite{tinyllm} evaluates sub-3B models on function calling, finding significant degradation with large tool sets. The SLM Survey~\cite{slm_survey} notes models under 20B are ``sufficient for schema-constrained accuracy'' but struggle with large registries.

\textbf{Gap.} No existing work addresses \emph{adaptive tool presentation} based on model capability. Prior approaches fine-tune models for better tool calling (Gorilla~\cite{gorilla}, ToolLLM~\cite{toollm}). We adapt the presentation instead --- zero training cost, works with any model.

% ======================================================================
\section{Method}

\subsection{Model Capability Tiers}

We classify models by parameter count into four tiers (Table~\ref{tab:tiers}). Tier detection parses parameter count from model name strings (e.g., \texttt{qwen3.5:9b} $\rightarrow$ 9B $\rightarrow$ Medium).

\begin{table}[h]
\centering
\caption{Model capability tiers.}
\label{tab:tiers}
\small
\begin{tabular}{@{}lll@{}}
\toprule
Tier & Parameters & Key constraint \\
\midrule
Tiny & $<$ 1.5B & Struggles with $>$10 tools \\
Small & 1.5--4B & Limited structured JSON \\
Medium & 4--14B & Reliable with 8--25 tools \\
Large & 14B+ & Handles full tool sets \\
\bottomrule
\end{tabular}
\end{table}

\subsection{Tool Families}

80 tools are organized into 8 semantic families: filesystem~(10), code and git~(12), web and API~(10), data~(10), communication~(10), system~(10), DevOps~(10), AI and memory~(8).

\subsection{Presentation Strategies}

We evaluate seven strategies, organized by mechanism:

\textbf{Filtering strategies} (reduce tool count):
\begin{itemize}[nosep,leftmargin=*]
  \item \textbf{Baseline.} All 80 tools via native API.
  \item \textbf{Semantic top-8.} YantrikDB~\cite{yantrikdb} vector similarity ranks tools; top-8 presented.
  \item \textbf{Semantic top-4.} Same as above, $K{=}4$.
\end{itemize}

\textbf{Routing strategies} (select subset by category):
\begin{itemize}[nosep,leftmargin=*]
  \item \textbf{Family oracle.} Correct family provided (upper bound).
  \item \textbf{Family detected.} Auto-detected via YantrikDB similarity.
\end{itemize}

\textbf{Presentation strategies} (change appearance, not set):
\begin{itemize}[nosep,leftmargin=*]
  \item \textbf{Semantic reorder.} All 80 tools, sorted by similarity.
  \item \textbf{Reorder + category hint.} Reorder + system prompt hint.
  \item \textbf{Hybrid (8+72).} Top-8 get full descriptions; 72 appear as name-only entries.
\end{itemize}

\subsection{Semantic Ranking via YantrikDB}

Tool descriptions are pre-embedded into YantrikDB's HNSW vector index using all-MiniLM-L6-v2 embeddings. At query time, the user prompt is embedded and tools are ranked by cosine similarity. This ranking drives the semantic strategies.

% ======================================================================
\section{Experimental Setup}

\subsection{Models}

\begin{table}[h]
\centering
\caption{Evaluated models.}
\label{tab:models}
\small
\begin{tabular}{@{}llll@{}}
\toprule
Model & Params & Tier & Host \\
\midrule
qwen2.5:1.5b & 1.5B & Tiny & local \\
qwen3.5:9b & 9B & Medium & remote \\
gpt-oss:20b & 20B & Large & local \\
qwen3.5:35b & 35B & Large & remote \\
\bottomrule
\end{tabular}
\end{table}

\subsection{Evaluation Protocol}

\begin{itemize}[nosep,leftmargin=*]
  \item \textbf{API}: Ollama \texttt{/api/chat} with \texttt{tools} parameter (native tool calling)
  \item \textbf{Temperature}: 0.0 (deterministic)
  \item \textbf{Prompts}: 50 natural-language requests spanning all 8 families
  \item \textbf{Metrics}: tool selection accuracy (exact match), family accuracy, prompt tokens, latency
  \item \textbf{Total}: 1,000+ inference calls across three rounds
\end{itemize}

\subsection{Methodological Note}

We use Ollama's \textbf{native tool calling API}, not text-injected descriptions. Early experiments using text injection (\texttt{/api/generate}) produced 10--20pp different results, confirming that benchmark methodology must match production usage. Models are trained for structured tool calling; text injection does not reflect deployment behavior.

% ======================================================================
\section{Results}

\subsection{Main Results}

Table~\ref{tab:accuracy_filter} shows filtering and routing strategy accuracy. Table~\ref{tab:accuracy_present} shows presentation strategy results.

\begin{table}[h]
\centering
\caption{Filtering \& routing strategies --- accuracy (\%).}
\label{tab:accuracy_filter}
\small
\begin{tabular}{@{}llrrrrr@{}}
\toprule
Model & Tier & Base & S-8 & S-4 & Oracle & Det. \\
\midrule
1.5B & Tiny & 50 & \textbf{64} & 64 & \textbf{70} & 54 \\
9B & Med & 80 & 72 & 72 & \textbf{86} & 64 \\
20B & Lrg & 80 & 70 & 70 & \textbf{84} & 58 \\
35B & Lrg & \textbf{88} & 76 & 78 & \textbf{88} & 64 \\
\bottomrule
\end{tabular}
\end{table}

\begin{table}[h]
\centering
\caption{Presentation strategies --- accuracy (\%, 1.5B \& 20B).}
\label{tab:accuracy_present}
\small
\begin{tabular}{@{}lrrrr@{}}
\toprule
Model & Base & Reorder & +Hint & Hybrid \\
\midrule
1.5B & 50 & 54 & 54 & \textbf{60} \\
20B & 80 & 84 & \textbf{88} & 76 \\
\bottomrule
\end{tabular}
\end{table}

\begin{table}[h]
\centering
\caption{Average prompt tokens by strategy.}
\label{tab:tokens}
\small
\begin{tabular}{@{}lrrrr@{}}
\toprule
Model & Base & S-8 & Oracle & Hybrid \\
\midrule
1.5B & 3,408 & 444 & 540 & $\sim$1,800 \\
9B & 5,272 & 724 & 872 & --- \\
20B & 2,143 & 310 & 368 & $\sim$1,200 \\
35B & 5,272 & 724 & 872 & --- \\
\bottomrule
\end{tabular}
\end{table}

\subsection{The Decomposition: Family Routing Is the Bottleneck}

Table~\ref{tab:decomposition} presents the accuracy decomposition from Equation~\ref{eq:decomposition}.

\begin{table}[h]
\centering
\caption{Accuracy decomposition: $P(\text{tool}) = P(\text{fam}) \times P(\text{tool}|\text{fam})$.}
\label{tab:decomposition}
\small
\begin{tabular}{@{}llrrr@{}}
\toprule
Model & Strategy & $P$(fam) & $P$(tool) & $P$(t$|$f) \\
\midrule
1.5B & Baseline & 56\% & 50\% & \textbf{89\%} \\
1.5B & Semantic-8 & 70\% & 64\% & \textbf{91\%} \\
1.5B & Oracle & 74\% & 70\% & \textbf{95\%} \\
\midrule
9B & Baseline & 82\% & 80\% & \textbf{98\%} \\
9B & Oracle & 88\% & 86\% & \textbf{98\%} \\
\midrule
20B & Baseline & 84\% & 80\% & \textbf{95\%} \\
\midrule
35B & Baseline & 90\% & 88\% & \textbf{98\%} \\
\bottomrule
\end{tabular}
\end{table}

\textbf{Finding~1: Within-family accuracy is consistently high (89--98\%).} Even the 1.5B model selects the correct tool 89\% of the time when presented with the correct family. The gap between model sizes is primarily in family routing (56\% vs 90\%), not in tool selection given the right family.

\textbf{Finding~2: Family routing is the dominant error source.} Error analysis for the 1.5B baseline shows: 22/25 errors are wrong-family, only 3/25 are right-family-wrong-tool, plus 20 prompts received no tool call. The selection mechanism works; the search space is the problem.

\subsection{Strategy Effectiveness by Model Scale}

\textbf{Finding~3: Different strategies win at different scales.}
\begin{itemize}[nosep,leftmargin=*]
  \item \textbf{Tiny (1.5B):} Hybrid wins (+10pp). Detailed descriptions for top-8 semantically relevant tools, name-only for the rest.
  \item \textbf{Large (20B):} Reorder + hint wins (+8pp). Semantic ordering + system prompt achieves 88\% --- matching 35B baseline.
  \item \textbf{No strategy universally dominates.} Hybrid hurts 20B ($-$4pp); reorder+hint doesn't help 1.5B (+4pp only).
\end{itemize}

\textbf{Finding~4: Semantic filtering helps small models but can hurt large ones.} Semantic top-8 improves 1.5B from 50\% to 64\% (+14pp) but drops 20B from 80\% to 70\% ($-$10pp). Aggressive filtering removes tools the larger model could have found. The hybrid approach solves this by keeping all tools accessible.

\subsection{Token Efficiency}

\textbf{Finding~5: Token savings of 83--92\% are achievable.} Filtering strategies reduce tokens by 83--92\%. Hybrid achieves $\sim$47\% reduction. Reorder-based strategies use baseline tokens but improve accuracy --- a pure win.

\subsection{Automated Family Detection}

\textbf{Finding~6: Automated detection underperforms baseline.} YantrikDB-based family detection achieves only 54--64\% accuracy --- worse than baseline everywhere. When detection is wrong, the correct tool is absent, causing hard failure. The oracle ceiling (+4--20pp) shows the potential; bridging this gap is the open challenge.

% ======================================================================
\section{Production: YantrikOS}

\subsection{System Overview}

YantrikOS is an AI-native desktop operating system (Rust, single binary) with 116+ tools across 48 categories, supporting models from 0.8B to 35B+. The \texttt{ModelCapabilityProfile} structure adapts six dimensions: max tools per prompt, tool call format (MCQ/JSON/native), slot extraction mode, family routing, context budget, and confidence thresholds.

\subsection{Resolving the Detection Bottleneck}

YantrikOS resolves family detection through \texttt{discover\_tools} --- a meta-tool enabling iterative navigation:
\begin{enumerate}[nosep,leftmargin=*]
  \item \texttt{discover\_tools()} $\rightarrow$ category summary
  \item \texttt{discover\_tools(category="filesystem")} $\rightarrow$ family tools
  \item Model selects and invokes
\end{enumerate}

This multi-turn pattern allows self-correction when the initial choice is wrong. Our single-shot benchmark cannot capture this, explaining the gap between automated detection (54--64\%) and production effectiveness.

\subsection{Model Family Awareness}

YantrikOS implements per-family chat templates (Qwen, Llama, Nemotron, Gemma, Phi) ensuring tool call format matches model training --- orthogonal to tier-based routing.

% ======================================================================
\section{SDK for Tier-Aware Tools}

The Yantrikos SDK enforces tier-aware tool design:

\begin{verbatim}
from yantrikos import BaseTool, Tier

class FileReadTool(BaseTool):
    name = "file_read"
    descriptions = {
        Tier.S: "Read file",
        Tier.L: "Read file with encoding
                 and line number control",
    }
    parameters = {
        Tier.S: {"path": str},
        Tier.L: {"path": str,
                 "encoding": str,
                 "line_numbers": bool},
    }
\end{verbatim}

\textbf{Guidelines:} (1)~Provide at least two description lengths. (2)~Stratify parameters by tier. (3)~Keep families semantically distinct. (4)~Use native tool calling APIs.

% ======================================================================
\section{Limitations}

\begin{enumerate}[nosep,leftmargin=*]
  \item \textbf{Single-shot evaluation.} Iterative discovery (multi-turn) likely achieves higher accuracy.
  \item \textbf{Four models, two families.} Results may differ for Llama, Gemma, Phi.
  \item \textbf{English only.} Cross-lingual evaluation needed.
  \item \textbf{General-purpose embeddings.} Specialized tool-routing embeddings could improve detection.
  \item \textbf{Fixed tool registry.} Dynamic environments may need different strategies.
  \item \textbf{No confidence intervals.} With 50 prompts, reported accuracies have $\pm$7pp margin at 95\% CI.
\end{enumerate}

% ======================================================================
\section{Threats to Validity}

\textbf{Prompt sensitivity.} Tool selection may be sensitive to exact prompt wording. Our 50 prompts are hand-crafted and may not represent the full distribution of real user requests.

\textbf{Model-provider dependence.} Results are specific to Ollama's native tool calling implementation. Other providers (vLLM, TGI, cloud APIs) may produce different tool calling behavior.

\textbf{Semantic ranker quality.} YantrikDB with all-MiniLM-L6-v2 provides general-purpose embeddings. Tool-specific embedding models could change the results substantially.

\textbf{Tool registry design.} Our 80-tool registry was designed to be realistic but is not drawn from a production deployment. Tool description quality and family boundary clarity directly affect results.

% ======================================================================
\section{Conclusion}

We presented Tier-Based Adaptive Tool Routing, evaluated across 1,000+ inference calls with native tool calling. Our decomposition reveals that tool selection is primarily a family routing problem --- even 1.5B models achieve 89\% accuracy within the correct family. Key results:

\begin{enumerate}[nosep,leftmargin=*]
  \item \textbf{Hybrid presentation} (8~detailed + 72~name-only) improves 1.5B accuracy from 50\% to 60\%.
  \item \textbf{Semantic reorder + hint} improves 20B from 80\% to 88\%, matching 35B performance.
  \item \textbf{No single strategy dominates} --- optimal presentation is scale-dependent.
  \item \textbf{Token savings of 83--92\%} are universal across filtering strategies.
  \item \textbf{Automated family detection remains the open challenge} at 54--64\%, solvable through iterative LLM-driven discovery as demonstrated in YantrikOS.
\end{enumerate}

The framework requires no fine-tuning and works with any model supporting tool calling. Code, data, and SDK: \url{https://github.com/yantrikos/tier}

% ======================================================================
\section*{Acknowledgments}

This work was conducted independently. All models were evaluated on consumer hardware (Apple Silicon MacBook Pro + homelab with Ollama). The author thanks the open-source communities behind Qwen, Ollama, and Hugging Face for making local LLM experimentation accessible.

% ======================================================================
\section*{Reproducibility Statement}

\textbf{Artifacts.} We release: (i)~the benchmark harness (\texttt{harness\_v3.py}), (ii)~the full tool registry (80~tools), (iii)~all 50 test prompts, (iv)~raw results as JSONL (1,000+ data points), (v)~the Yantrikos SDK.

\textbf{Environment.} Local: Apple Silicon, Ollama v0.6+. Remote: homelab server, Ollama v0.6+. Models served via \texttt{/api/chat} with native tool calling. Temperature: 0.0. Max generation: 256 tokens.

\textbf{Procedure.} To replicate:
\begin{verbatim}
git clone https://github.com/yantrikos/tier
cd tier
pip install yantrikdb sentence-transformers
python benchmarks/harness_v3.py
\end{verbatim}

% ======================================================================
\begin{thebibliography}{10}

\bibitem{agentflux}
R.~Kadekodi et~al., ``AgentFlux: Decoupled Fine-Tuning and Inference for On-Device Agentic Systems,'' \textit{arXiv:2510.00229}, 2025.

\bibitem{toollm}
Y.~Qin et~al., ``ToolLLM: Facilitating Large Language Models to Master 16000+ Real-world APIs,'' \textit{arXiv:2307.16789}, 2023.

\bibitem{gorilla}
S.~Patil et~al., ``Gorilla: Large Language Model Connected with Massive APIs,'' \textit{arXiv:2305.15334}, 2023.

\bibitem{qwen25}
Qwen Team, ``Qwen2.5 Technical Report,'' \textit{arXiv:2412.15115}, 2024.

\bibitem{llama3}
Meta AI, ``Llama 3 Model Card,'' 2024.

\bibitem{tinyllm}
``TinyLLM: Evaluation and Optimization of Small Language Models for Agentic Tasks on Edge Devices,'' \textit{arXiv:2511.22138}, 2025.

\bibitem{slm_survey}
``Small Language Models for Agentic Systems: A Survey,'' \textit{arXiv:2510.03847}, 2025.

\bibitem{yantrikdb}
P.~Sarkar, ``YantrikDB: A Cognitive Memory Engine for Persistent AI Systems,'' Zenodo, DOI:~10.5281/zenodo.18793952, 2026. U.S. Patent Application 19/573,392.

\bibitem{openai_fc}
OpenAI, ``Function Calling and Structured Outputs,'' OpenAI API Documentation, 2024.

\bibitem{sdf}
P.~Sarkar, ``Convert Once, Consume Many: SDF for Cacheable, Typed Semantic Extraction from Web Pages,'' Zenodo, DOI:~10.5281/zenodo.18559223, 2026.

\end{thebibliography}

% ======================================================================
\appendix

\section{Error Taxonomy (1.5B Baseline)}

\begin{table}[h]
\centering
\small
\begin{tabular}{@{}lr@{}}
\toprule
Error Type & Count (/50) \\
\midrule
Correct tool & 25 \\
Wrong family & 22 \\
Right family, wrong tool & 3 \\
No tool call & 20 \\
\bottomrule
\end{tabular}
\end{table}

88\% of errors (22/25) are wrong-family errors. Only 12\% (3/25) are within-family disambiguation failures. This confirms that family routing, not tool selection, is the dominant bottleneck.

\section{Six Adaptation Dimensions (YantrikOS)}

\begin{table}[h]
\centering
\small
\begin{tabular}{@{}lllll@{}}
\toprule
Dimension & Tiny & Small & Med & Large \\
\midrule
Max tools & 3--5 & 5--10 & 8--25 & 50+ \\
Call format & MCQ & JSON & Native & Native \\
Slot mode & KV & JSON & JSON & JSON \\
Family route & Yes & Yes & Yes & Opt. \\
Context & 512t & 1024t & 2048t & 4096t \\
Confidence & 0.95 & 0.85 & 0.80 & 0.70 \\
\bottomrule
\end{tabular}
\end{table}

\end{document}
